\documentclass[french]{book}
\usepackage{teach}
\usepackage{verbatim}
\usepackage{amsmath,amsfonts,amssymb}
\usepackage{pgfplots}
\RequirePackage{graphicx}
\RequirePackage{ifpdf}
 \ifpdf
   \DeclareGraphicsRule{*}{mps}{*}{}
 \fi
\usepackage{fancyhdr}
\usepackage{pstricks}
\usepackage{amsmath}
\usepackage{amssymb}
\usepackage{amsfonts}
\usepackage{amsthm}

\usepackage{tkz-tab}
\usepackage{mathrsfs}
\usepackage{array}
\usepackage{multicol}
\setlength{\multicolsep}{3pt}
\usepackage{graphicx}
\graphicspath{{Figures/}}
\usepackage{pstricks,pst-plot,pst-text,pst-tree,pst-eucl}
\usepackage[all]{xy}
\usepackage{fancybox}
\usepackage{color}
\usepackage{hyperref}
%  \usepackage[frenchb]{babel}
\usepackage{subfig}
\pgfplotsset{compat=newest}
%\usepackage{geometry}
%\geometry{top=1cm, bottom=1cm, left=2cm, right=2cm}
\usepackage[latin1]{inputenc}
\usepackage[french]{babel}
\redefinecolor{thm}{title}{bgcolor}{alizarin}
\redefinecolor{prop}{title}{bgcolor}{meatbrown}
\redefinecolor{defn}{title}{bgcolor}{olivine}
\definecolor{alizarin}{rgb}{0.82, 0.1, 0.26}
\definecolor{meatbrown}{rgb}{0.9, 0.72, 0.23}
\definecolor{olivine}{rgb}{0.6, 0.73, 0.45}
\usepackage{mathrsfs}
%%
\usepackage{yhmath} 
\newcommand{\R}{\mathbb {R}}
%\newcommand{\C}{\mathds {C}}
\newcommand{\N}{\mathbb {N}}
\newcommand{\Z}{\mathbb {Z}}
\newcommand{\Q}{\mathbb {Q}}
\newcommand{\D}{\mathbb {D}}
%%
\newcommand{\se}{\geq}
\newcommand{\ie}{\leq}
\newcommand{\tm}{\times}

%\newcommand{\ell}{l}
%\newcommand{\ell'}{l'}
%%

\newcommand{\sysd}[2]{%
       \left\{
       \begin{aligned}
          #1\\
          #2\\
       \end{aligned}
       \right.%
       }

\newcommand{\ds}{\displaystyle}
\newcommand{\limn}{\lim_{n\rightarrow +\infty}}
\newcommand{\limo}[1][x]{\ds\lim_{#1\rightarrow 0}}
\newcommand{\limite}[2][x]{\ds\lim_{#1\rightarrow #2}}
\newcommand{\limpinf}[1][x]{\ds\lim_{#1\rightarrow +\infty}}
\newcommand{\limminf}[1][x]{\ds\lim_{#1\rightarrow -\infty}}
\newcommand{\limiteg}[2][x]{\ds\lim_{#1\xrightarrow{<}#2}}
\newcommand{\limited}[2][x]{\ds\lim_{#1\xrightarrow{>}#2}}
%%
\newcommand{\vect}[1]{\overrightarrow{#1}}
\newcommand{\arc}[1]{\wideparen{#1}}
\newcommand{\oij}{(\textit{O},{\overrightarrow{i}},{\overrightarrow{j}})}
%\newcommand{\imath}{\overrightarrow{i}}
%\newcommand{\jmath}{\overrightarrow{j}}
%
\newcommand{\cp}[2]{%
        \begin{pmatrix}
        #1\\
        #2
        \end{pmatrix}%
        }
%
\newcommand{\calb}{\mathscr{B}}
\newcommand{\calc}{\mathscr{C}}
\newcommand{\cald}{\mathscr{D}}
\newcommand{\cale}{\mathscr{E}}
\newcommand{\calf}{\mathscr{F}}
\newcommand{\calp}{\mathscr{P}}
\newcommand{\cals}{\mathscr{S}}
\newcommand{\calt}{\mathscr{T}}
%
\newcommand{\suite}[1][u]{\left( #1_{n} \right)_{n\in\N}}
\newcommand{\suitep}[1][u]{\left( #1_{n} \right)_{n\in\N^{*}}}
\usetikzlibrary{arrows}

\newcommand{\poly}[1][a]{#1_0+#1_1x+\cdots+#1_nx^n}

\begin{document}
\tableofcontents
\chapter{Trigonom\'etrie}
\section{Lignes trigonom\'etriques}

\subsection{Cosinus et Sinus}
\subsubsection*{Rappels}
\begin{minipage}{10cm}
Soit $\oij$ un rep\`ere orthonormal direct (c'est-�-dire tel que $(\vect{\imath},\vect{\jmath})=\frac{\pi}{2}+2k\pi$), $\calc$ le cercle trigonom\'etrique de centre $O$.
\par $A$ et $B$ sont les points tels que $\vect{OA}=\vect{\imath}$ et $\vect{OB}=\vect{\jmath}$.

\vspace{1ex}
\begin{defn}
Pour tout r\'eel $x$, il existe un unique point $M$ de $\calc$ tel que $x$ soit une mesure de l'angle $(\vect{\imath},\vect{OM})$.
\par $\cos x$ est l'abscisse de $M$ dans $\oij$.
\par $\sin x$ est l'ordonn\'ee de $M$ dans $\oij$.
\end{defn}
\end{minipage}
\hfill
\begin{minipage}{4cm}
\includegraphics[scale=1.3]{CoursTrigoFigures.1}
\end{minipage}

\subsubsection*{Propri\'et\'es imm\'ediates}
\begin{prop}
Pour tout $x\in\R$ :

\begin{tabular}{*{3}{m{0.3\linewidth}}}
$-1\ie \cos x \ie 1$ \par $-1\ie \sin x \ie 1$ & $\cos(x+2k\pi)=\cos x$ \par $\sin(x+2k\pi)=\sin x$ & $\cos^2 x+\sin^2x=1$
\end{tabular}
\end{prop}

\begin{demo}
$M$ appartient au cercle $\calc$ dont une \'equation cart\'esienne est $x^2 + y^2 = 1$. Les coordonn\'ees de M sont $(\cos x;\sin x)$ et v\'erifient l?\'equation de $\calc$ d'o\`u le r\'esultat.
\end{demo}



\subsubsection*{Cosinus et sinus d'un angle orient\'e de vecteurs}
Soit $(\vect{u},\vect{v})$ un angle orient\'e et $x$ une de ses mesures. Les autres mesures de $(\vect{u},\vect{v})$ sont donc de la forme $x+2k\pi$. Or $\cos(x+2k\pi)=\cos x$ et $\sin(x+2k\pi)=\sin x$. On a donc la d\'efinition suivante :

\vspace{1ex}
\begin{defn}
Le cosinus (resp. le sinus) d'un angle orient\'e de vecteurs est le cosinus (resp. le sinus) d'une quelconque de ses mesures.
\end{defn}

\newpage%%%%%%%%%%%%%%%%%%%%%%%%%%%%%%%%%%%%%%%%%%%%%%%%%%%%%%%%%%%%%%%%%%%%%%%%%%%%%%%%%%%%%%%%%%%%%%%%%%%%
\subsection{Angles associ\'es}
L'utilisation de sym\'etries dans le cercle trigonom\'etrique permet d'\'etablir :% les relations suivantes :

\vspace{1ex}
\begin{prop}
Pour tout r\'eel $x$,

\vspace{1ex}
$\sysd{\cos(-x)=\cos x}{\sin(-x)=-\sin x}$ \hfill $\sysd{\cos(x+\pi)=-\cos x}{\sin(x+\pi)=-\sin x}$ \hfill $\sysd{\cos(\pi-x)=-\cos x}{\sin(\pi-x)=\sin x}$

\hfill $\sysd{\cos\left(\dfrac{\pi}{2}-x\right)=\sin x}{\sin\left(\dfrac{\pi}{2}-x\right)=\cos x}$ \hfill $\sysd{\cos\left(\dfrac{\pi}{2}+x\right)=-\sin x}{\sin\left(\dfrac{\pi}{2}+x\right)=\cos x}$ \hfill\phantom{.}

\end{prop}

\includegraphics{CoursTrigoFigures.2} \hfill \includegraphics{CoursTrigoFigures.3}

\subsection{Valeurs particuli\`eres}
� l'aide de consid\'erations g\'eom\'etriques, on peut obtenir les valeurs suivantes :

\renewcommand{\arraystretch}{3}
\begin{center}
\begin{tabular}{|c|c|c|c|c|c|c|} \hline
$x$ & 0 & $\dfrac{\pi}{6}$ & $\dfrac{\pi}{4}$ & $\dfrac{\pi}{3}$ & $\dfrac{\pi}{2}$ & $\pi$ \\ \hline
$\cos(x)$ & 1 & $\dfrac{\sqrt{3}}{2}$ & $\dfrac{\sqrt{2}}{2}$ & $\dfrac{1}{2}$ & 0 & -1 \\ \hline
$\sin(x)$ & 0 & $\dfrac{1}{2}$ & $\dfrac{\sqrt{2}}{2}$ & $\dfrac{\sqrt{3}}{2}$ & 1 & 0 \\ \hline
\end{tabular}
\end{center}
\renewcommand{\arraystretch}{1}
\subsection{Formules d'addition}
\begin{prop}
Quels que soient les nombres r\'eels $a$ et $b$,\\

\vspace{-1ex}
\begin{minipage}{0.48\linewidth}
\begin{equation}
\cos(a-b)=\cos a \cos b + \sin a \sin b
\end{equation}
\begin{equation}
\cos(a+b)=\cos a \cos b - \sin a \sin b
\end{equation}
\end{minipage}
\begin{minipage}{0.48\linewidth}
\begin{equation}
\sin(a-b)=\sin a \cos b - \cos a \sin b
\end{equation}
\begin{equation}
\sin(a+b)=\sin a \cos b + \cos a \sin b
\end{equation}
\end{minipage}
\end{prop}

\begin{demo}
\underline{\'equation (1.1)}: Soit $\oij$ un rep\`ere orthonormal, $\calc$ le cercle trigonom\'etrique de centre $O$ et $A$ et $B$ les points de $\calc$ tels que $(\vect{\imath},\vect{OA})=a$ et $(\vect{\imath},\vect{OB})=b$. 

\par Calculons de deux fa�ons diff\'erentes le produit scalaire $\vect{OA}.\vect{OB}$\\

\begin{itemize}
\item $\vect{OA}.\vect{OB}=OA\times OB\times\cos(\vect{OA},\vect{OB})$
\par Or $(\vect{OA},\vect{OB})=(\vect{OA},\vect{\imath})+(\vect{\imath},\vect{OB})+2k\pi=-a+b+2k\pi$ et $OA=OB=1$
\par Ainsi $\vect{OA}.\vect{OB}=1\times1\times\cos(b-a)=\cos(a-b)$
\item Les coordonn\'ees de $\vect{OA}$ sont $\cp{\cos a}{\sin a}$ et les coordonn\'ees de $\vect{OB}$ sont $\cp{\cos b}{\sin b}$
\par On a donc : $\vect{OA}.\vect{OB}=\cos a \cos b + \sin a \sin b$
\end{itemize}

\underline{Conclusion }: $\cos(a-b)=\cos a \cos b + \sin a \sin b$\\

\underline{\'equation (1.2)}: $\cos(a+b)=\cos(a-(-b))=\cos a \cos (-b) + \sin a \sin (-b)=\cos a \cos b - \sin a \sin b$\\

\underline{\'equation (1.3) }: $\sin(a-b)=\cos\left( \dfrac{\pi}{2}-(a-b) \right)=\cos\left( \left(\dfrac{\pi}{2}-a \right)+b\right) = \cos\left(\dfrac{\pi}{2}-a \right)\cos b-\sin\left(\dfrac{\pi}{2}-a \right)\sin b = \sin a \cos b - \cos a \sin b$ \\

\underline{\'equation (1.4) }: $\sin(a+b)=\sin(a-(-b))=\sin a \cos (-b) - \cos a \sin (-b) = \sin a \cos b + \cos a \sin b$
\end{demo}


%%%%%%%%%%%%%%%%%%AJOUT AVRIL 2008%%%%%%%%%%%%%%%%%%%%%%%%%%%%%%%%%%%%%%%%%
\subsection{Formules de duplication}
\begin{prop}
Quel que soit le r\'eel $x$,
\begin{align}
 \cos(2x)&=\cos^2x-\sin^2x=2\cos^2x-1=1-2\sin^2x\label{cos2x} \\
 \sin(2x)&=2\sin x\cos x\label{sin2x}
\end{align}
\end{prop}

\begin{demo}
(\ref{cos2x}) : Quel que soit le r\'eel $x$,
 \[\cos(2x)=\cos(x+x)=\cos x\cos x-\sin x\sin x=\cos^2x-\sin^2x\]
En utilisant les relations $\sin^2x=1-\cos^2x$ et $\cos^2x=1-\sin^2x$, on obtient successivement :
\[\cos(2x)=2\cos^2x-1\quad\text{puis}\quad\cos(2x)=1-2\sin^2x\]
(\ref{sin2x}) : Quel que soit le r\'eel $x$,
 \[\sin(2x)=\sin(x+x)=\sin x\cos x+\cos x\sin x=2\sin x\cos x\]

\end{demo}

\section{R\'esolution d'\'equations}
\begin{prop}
Quels que soient les r\'eels $a$ et $b$,

%\begin{equation*}
 $\cos a=\cos b \Leftrightarrow
   \left|\text{ou }
     \begin{array}{@{}l}
	 a=b+k\times2\pi \\
	 a=-b+k\times2\pi
     \end{array}
   \right.$
     \hfill et\hfill
 $\sin a=\sin b \Leftrightarrow
   \left|\text{ou }
     \begin{array}{@{}l}
	 a=b+k\times2\pi \\
	 a=\pi-b+k\times2\pi
     \end{array}
   \right.$
%\end{equation*}
\end{prop}

\begin{exemple}
R\'esoudre dans $]-\pi~;~\pi[$ l'\'equation $\sin\left(x-\frac{\pi}{3}\right)=\frac{1}{2}$
\end{exemple}
\begin{align*}
\sin\left(x+\frac{\pi}{3}\right)=\frac{1}{2} &\Leftrightarrow \sin\left(x-\frac{\pi}{3}\right)=\sin \frac{\pi}{6}\\
&\Leftrightarrow
    \left|\text{ou }
     \begin{array}{@{}l}
	 x-\frac{\pi}{3}=\frac{\pi}{6}+k\times2\pi \\
	 x-\frac{\pi}{3}=\pi-\frac{\pi}{6}+k\times2\pi
     \end{array}
   \right.\\
&\Leftrightarrow
    \left|\text{ou }
     \begin{array}{@{}l}
	 x=\frac{\pi}{2}+k\times2\pi \\
	 x=\frac{7\pi}{6}+k\times2\pi
     \end{array}
   \right.
\end{align*}
Sur l'intervalle $]-\pi~;~\pi[$, on a $\cals=\{\frac{\pi}{2};-\frac{5\pi}{6}\}$

%%%%%%%%%%%%%%%%%%FIN AJOUT%%%%%%%%%%%%%%%%%%%%%%%%%%%%%%%%%%%%%%%%%%%%%%%%%

\newpage
\section{Rep\'erage polaire}
\subsection{Coordonn\'ees polaires d'un point}
\begin{minipage}{11cm}
Soit $\oij$ un rep\`ere orthonormal direct. Si $M$ est un point distinct du point $O$ alors $M$ peut-\^etre rep\'er\'e par l'angle $\theta=(\vect{\imath},\vect{OM})$ et la longueur $r=OM$.
\par R\'eciproquement, la donn\'ee d'un couple $(r;\theta)$ avec $r>0$ et $\theta\in\R$ d\'etermine un seul point $M$ tel que $\theta=(\vect{\imath},\vect{OM})$ et $r=OM$.
\end{minipage}
\hfill
\begin{minipage}{3.2cm}
\includegraphics[scale=1.3]{CoursTrigoFigures.4}
\end{minipage}

%\vspace{1ex}
\begin{defn}
Pour tout point $M$ distinct de $O$, un couple $(r;\theta)$ tel que $\theta=(\vect{\imath},\vect{OM})$ et $r=OM$ est appel\'e couple de coordonn\'ees polaires de $M$ dans le rep\`ere polaire $(O,\vect{\imath})$.
\end{defn}


\subsection{Lien entre coordonn\'ees cart\'esiennes et polaires}
Soit $\oij$ un rep\`ere orthonormal direct.

\begin{prop}
Si $M$ est un point ayant pour coordonn\'ees cart\'esiennes $(x;y)$ dans le rep\`ere $\oij$ et pour coordonn\'ees polaires $(r;\theta)$ alors :
\begin{center}
$r=\sqrt{x^{2}+y^{2}}$ \qquad  ; \qquad $x=r\cos \theta$ \qquad ; \qquad $y=r\sin \theta$
\end{center}
\end{prop}

\begin{demo}
Notons $\calc$ le cercle trigonom\'etrique de centre $O$. La demi-droite $[OM)$ coupe $\calc$ en $N$. On a donc $\vect{OM}=r\vect{ON}$.
\par $N\in\calc$ donc ses coordonn\'ees cart\'esiennes sont $(\cos \theta ; \sin \theta)$. Celles de $M$ sont donc $(r\cos \theta ; r\sin \theta)$.
\par Par unicit\'e des coordonn�es, $x=r\cos \theta$ et $y=r\sin \theta$.

De plus $OM^{2}=x^{2}+y^{2}$ et $OM=r$ donc $r^{2}=x^{2}+y^{2}$.
\end{demo}


\section{Etude de fonctions trigonom�triques }
\subsection{Ensemble de d�finition, p�riodicit� et parit� }
\begin{prop}
Les fonction sinus et cosinus sont d�finies, continues, $2\pi$-p�riodiques et d�rivables sur $\mathbb{R}$.
\begin{itemize}
\item[$\bullet$] La fonction cosinus est paire. Sa courbe est donc sym�trique par rapport � l?axe des ordonn�es en rep�re orthogonal.
\item[$\bullet$] La fonction sinus est impaire. Sa courbe est donc sym�trique par rapport � l?origine du repe?re.
 \end{itemize}
\end{prop}
\subsubsection{Courbes repr�sentatives}
\begin{center}
\includegraphics{../Figures/trigo.1}\\
\vspace{1cm}
\includegraphics{../Figures/trigo.2}
\end{center}
\subsection{D�riv�es et variation de fonctions trigonom�triques  }

\begin{thm}
Pour tout $x \in \mathbb{R}$, $\sin'(x)=cos(x)$ $\cos'(x)=- \sin(x)$
\end{thm}

%\begin{tikzpicture}
  % \tkzTabInit{$x$ /1, $-\sin(x)$ /1, $\cos(x)$ /1.5} {$-\pi$ , $0$, $\pi$}
  % \tkzTabLine{, +, 0, -, } 
  % \tkzTabVar{+/ 1, 1, -/ -1}
 %  \tkzTabIma{1}{3}{2}{$1$}
%\end{tikzpicture}

\subsubsection{Tableau de variation}
\begin{center}
\begin{tikzpicture}
   \tkzTabInit{$x$ / 1 , $\cos(x)$ / 1, $\sin(x)$ / 1.5}{$-\pi$, $-\dfrac{\pi}{2}$, 0,$\dfrac{\pi}{2}$, $\pi$}
   \tkzTabLine{, -, z,,+, , z,- }
   \tkzTabVar{+/ 0 ,-/ -1, R/ ,+/ 1 , -/ 0,}
   \tkzTabIma{2}{4}{3}{$0$}
\end{tikzpicture}

\vspace{5mm} 
 
\begin{tikzpicture}
   \tkzTabInit{$x$ / 1 , $- \sin(x)$ / 1, $\cos(x)$ / 1.5}{$-\pi$, $-\dfrac{\pi}{2}$, 0,$\dfrac{\pi}{2}$, $\pi$}
   \tkzTabLine{, +,,+ ,z, -, ,- }
   \tkzTabVar{-/ -1 , R/, +/ 1 ,R/  , -/ -1,}
 \tkzTabIma{1}{3}{2}{$0$}
  \tkzTabIma{3}{5}{4}{$0$}
\end{tikzpicture}

\end{center}

\end{document}
